\documentclass[10pt,a4paper]{amsart}
\usepackage{amsmath,amssymb,amsthm,graphicx}
% For commutative diagrams we use Xy-pic
\usepackage[all]{xy}

% We use 2cell for 2-commutative diagrams.
\xyoption{2cell}
% Theorem environments.
%
\theoremstyle{plain}
\newtheorem{theorem}[subsection]{Theorem}
\newtheorem{proposition}[subsection]{Proposition}
\newtheorem{lemma}[subsection]{Lemma}

\theoremstyle{definition}
\newtheorem{definition}[subsection]{Definition}
\newtheorem{example}[subsection]{Example}
\newtheorem{exercise}[subsection]{Exercise}
\newtheorem{situation}[subsection]{Situation}

\theoremstyle{remark}
\newtheorem{remark}[subsection]{Remark}
\newtheorem{remarks}[subsection]{Remarks}

\numberwithin{equation}{subsection}

% Macros
%
\def\lim{\mathop{\mathrm{lim}}\nolimits}
\def\colim{\mathop{\mathrm{colim}}\nolimits}
\def\Spec{\mathop{\mathrm{Spec}}}
\def\Hom{\mathop{\mathrm{Hom}}\nolimits}
\def\Ext{\mathop{\mathrm{Ext}}\nolimits}
\def\SheafHom{\mathop{\mathcal{H}\!\mathit{om}}\nolimits}
\def\SheafExt{\mathop{\mathcal{E}\!\mathit{xt}}\nolimits}
\def\Sch{\mathit{Sch}}
\def\Mor{\mathop{\mathrm{Mor}}\nolimits}
\def\Ob{\mathop{\mathrm{Ob}}\nolimits}
\def\Sh{\mathop{\mathit{Sh}}\nolimits}
\def\NL{\mathop{N\!L}\nolimits}
\def\CH{\mathop{\mathrm{CH}}\nolimits}
\def\proetale{{pro\text{-}\acute{e}tale}}
\def\etale{{\acute{e}tale}}
\def\QCoh{\mathit{QCoh}}
\def\Ker{\mathop{\mathrm{Ker}}}
\def\Im{\mathop{\mathrm{Im}}}
\def\Coker{\mathop{\mathrm{Coker}}}
\def\Coim{\mathop{\mathrm{Coim}}}

% Boxtimes
%
\DeclareMathSymbol{\boxtimes}{\mathbin}{AMSa}{"02}

%
% Macros for moduli stacks/spaces
%
\def\QCohstack{\mathcal{QC}\!\mathit{oh}}
\def\Cohstack{\mathcal{C}\!\mathit{oh}}
\def\Spacesstack{\mathcal{S}\!\mathit{paces}}
\def\Quotfunctor{\mathrm{Quot}}
\def\Hilbfunctor{\mathrm{Hilb}}
\def\Curvesstack{\mathcal{C}\!\mathit{urves}}
\def\Polarizedstack{\mathcal{P}\!\mathit{olarized}}
\def\Complexesstack{\mathcal{C}\!\mathit{omplexes}}
% \Pic is the operator that assigns to X its picard group, usage \Pic(X)
% \Picardstack_{X/B} denotes the Picard stack of X over B
% \Picardfunctor_{X/B} denotes the Picard functor of X over B
\def\Pic{\mathop{\mathrm{Pic}}\nolimits}
\def\Picardstack{\mathcal{P}\!\mathit{ic}}
\def\Picardfunctor{\mathrm{Pic}}
\def\Deformationcategory{\mathcal{D}\!\mathit{ef}}
\usepackage{cleveref}
\crefname{originallemma}{original lemma}{original lemmas}
\Crefname{originallemma}{Original lemma}{Original lemmas}
\makeatletter
\newlabel{more-algebra-lemma-application-formal-glueing}{{0ALK}{source}{}{sourcepointer-1}{}}
\newlabel{more-algebra-lemma-application-formal-glueing@cref}{{[originallemma][0][]0ALK}{[1][0][]source}{}{}{}}
\makeatother
\crefname{originalremark}{original remark}{original remarks}
\Crefname{originalremark}{Original remark}{Original remarks}
\makeatletter
\newlabel{more-algebra-remark-formal-glueing-algebras}{{05EU}{source}{}{sourcepointer-2}{}}
\newlabel{more-algebra-remark-formal-glueing-algebras@cref}{{[originalremark][0][]05EU}{[1][0][]source}{}{}{}}
\makeatother
\title{Gluing ring maps given at finitely many maximal ideals whose kernels and cokernels are killed locally}
\author{The Stacks Project Authors}
\date{stacks-project:chow.tex:local-snapshot-20261011; GNU FDL 1.2 or later}
\begin{document}
\maketitle
\noindent\textit{Editorial note.} All delivered mathematical text is the author's own, copied from the pinned Stacks Project chapter file chow.tex; no statement or proof step was rewritten, reordered or removed.\par\medskip
\noindent\textit{Editorial note.} The preamble is the chapter's own native preamble excerpt used for this source: the Xy-pic setup at lines 14-18 and the theorem declarations with the mathematical macros at lines 154-217 of sources/stacks-chow/preamble.tex. The quoted class line and the xr-hyper externaldocument block are not included because the assembler supplies a plain amsart wrapper; no macro was invented or redefined.\par\medskip
\noindent\textit{Editorial note.} The two cross-chapter references (a lemma and a remark of More on Algebra) are located in the project's own tag list staged in the snapshot (sources/stacks-chow/tags.txt), at the verbatim lines carrying their canonical tags.\par\medskip
\noindent\textit{Allowed standard prerequisites (editorial).}\begin{itemize}
\item Formal glueing of modules and algebras over a finite product of localizations. More on Algebra, Lemma more-algebra-lemma-application-formal-glueing (canonical Stacks tag 0ALK) and More on Algebra, Remark more-algebra-remark-formal-glueing-algebras (canonical Stacks tag 05EU); invoked at proof lines 548-555. Supply the existence of the A-module map and the algebra structure on the glued module; the reduction of the hypothesis to their input, i.e. the flatness of A -> A' with A/I = A'/IA', is inside the counted span.
\end{itemize}\medskip
\ifdefined\hypertarget\hypertarget{sourcepointer-1}{}\fi
\noindent\textit{Source reference.} Original lemma 0ALK: sources/stacks-chow/tags.txt, source lines 13373--13373. Cross-chapter standard prerequisite invoked once at proof lines 548-550 to glue the underlying module map; the verification that its hypotheses hold for A -> prod A\_\{m\_i\} is retained in the counted span.\par\medskip
\ifdefined\hypertarget\hypertarget{sourcepointer-2}{}\fi
\noindent\textit{Source reference.} Original remark 05EU: sources/stacks-chow/tags.txt, source lines 7003--7003. Cross-chapter standard prerequisite invoked once at proof lines 552-555 to endow the glued module with an A-algebra structure; it supplies the algebra structure only, while the reading off of finiteness, injectivity and surjectivity is retained in the counted span.\par\medskip
\begin{lemma}
\label{lemma-glue-at-max}
Let $A$ be a Noetherian ring. Let $\mathfrak m_1, \ldots, \mathfrak m_r$
be pairwise distinct maximal ideals of $A$. For $i = 1, \ldots, r$ let
$\varphi_i : A_{\mathfrak m_i} \to B_i$ be a ring map whose
kernel and cokernel are annihilated by a power
of $\mathfrak m_i$. Then there exists a ring map $\varphi : A \to B$ such
that
\begin{enumerate}
\item the localization of $\varphi$ at $\mathfrak m_i$ is
isomorphic to $\varphi_i$, and
\item $\Ker(\varphi)$ and $\Coker(\varphi)$ are annihilated
by a power of $\mathfrak m_1 \cap \ldots \cap \mathfrak m_r$.
\end{enumerate}
Moreover, if each $\varphi_i$ is finite, injective, or
surjective then so is $\varphi$.
\end{lemma}

\begin{proof}
Set $I = \mathfrak m_1 \cap \ldots \cap \mathfrak m_r$. Set
$A_i = A_{\mathfrak m_i}$ and $A' = \prod A_i$.
Then $IA' = \prod \mathfrak m_i A_i$ and $A \to A'$
is a flat ring map such that $A/I \cong A'/IA'$.
Thus we may use More on Algebra, Lemma
\ref{more-algebra-lemma-application-formal-glueing}
to see that there exists an $A$-module map $\varphi : A \to B$
with $\varphi_i$ isomorphic to the localization of $\varphi$
at $\mathfrak m_i$. Then we can use the discussion in
More on Algebra, Remark \ref{more-algebra-remark-formal-glueing-algebras}
to endow $B$ with an $A$-algebra structure
matching the given $A$-algebra structure on $B_i$.
The final statement of the lemma follows easily from
the fact that $\Ker(\varphi)_{\mathfrak m_i} \cong \Ker(\varphi_i)$
and $\Coker(\varphi)_{\mathfrak m_i} \cong \Coker(\varphi_i)$.
\end{proof}

\end{document}
